\SourceLeaf{311}
\[
\begin{aligned}
\log b={}&\log\beta+\log n+\tfrac18nn\lambda\lambda+\tfrac1{64}n^4\lambda^4+\tfrac1{384}n^6\lambda^6+\tfrac1{2048}n^8\lambda^8+\text{usw.}\\
&+\tfrac1{24}\beta\beta(nn+\tfrac14n^4\lambda\lambda+\tfrac1{16}n^6\lambda^4+\tfrac1{64}n^8\lambda^6+\text{usw.})\\
&+\tfrac{11}{2880}\beta^4(n^4+\tfrac12n^6\lambda\lambda+\tfrac3{16}n^8\lambda^4+\text{usw.})\\
&+\tfrac{191}{362880}\beta^6(n^6+\tfrac34n^8\lambda\lambda+\text{usw.})\\
&+\tfrac{2497}{29030400}\beta^8(n^8+\text{usw.})\\
&+\text{usw.}
\end{aligned}
\]
oder, indem man die Glieder gleicher Ordnung zusammenfasst,
\[
\begin{aligned}
\log b={}&\log\beta+\log n\\
&+\tfrac1{24}nn(\beta\beta+3\lambda\lambda)\\
&+\tfrac1{2880}n^4(11\beta^4+30\beta\beta\lambda\lambda+45\lambda^4)\\
&+\tfrac1{362880}n^6(191\beta^6+693\beta^4\lambda\lambda+945\beta\beta\lambda^4+945\lambda^6)\\
&+\tfrac1{29030400}n^8(2497\beta^8+11460\beta^6\lambda\lambda+20790\beta^4\lambda^4\\
&\qquad\qquad+18900\beta\beta\lambda^6+14175\lambda^8)\\
&+\text{usw.}
\end{aligned}
\]

\indent Um die Gleichungen I, II und III vollständig zu entwickeln, setzt man in ihnen noch ein
\[
\begin{aligned}
\log m&=\tfrac1{12}rr+\tfrac7{1440}r^4+\tfrac{31}{90720}r^6+\tfrac{127}{4838400}r^8+\text{usw.}\\
mm&=1+\tfrac16rr+\tfrac{17}{720}r^4+\tfrac{31}{10080}r^6+\text{usw.}\\
m^4&=1+\tfrac13rr+\tfrac3{40}r^4+\text{usw.}\\
m^6&=1+\tfrac12rr+\text{usw.}\\
&\text{usw.}\\
\log n&=-\tfrac1{24}rr-\tfrac1{2880}r^4-\tfrac1{181440}r^6-\tfrac1{9676800}r^8-\text{usw.}\\
nn&=1-\tfrac1{12}rr+\tfrac1{360}r^4-\tfrac1{20160}r^6+\text{usw.}\\
n^4&=1-\tfrac16rr+\tfrac1{80}r^4-\text{usw.}\\
n^6&=1-\tfrac14rr+\text{usw.}\\
&\text{usw.}
\end{aligned}
\]

\indent Für die Logarithmen von \(t,l,b\), genauer für die Unterschiede zwischen diesen Logarithmen und den Näherungswerten \(\log\tau,\log\lambda,\log\beta\), erhält man also zusammengesetzte Reihen, die folgendermaßen fortschreiten
\begin{list}{}{\setlength{\leftmargin}{1.4em}\setlength{\itemindent}{0pt}\setlength{\itemsep}{0pt}\setlength{\parsep}{0pt}}
\item für \(\log t\) nach den geraden Potenzen von \(\tau\) und \(r\) und deren Produkten,
\item für \(\log l\) ebenso nach den Potenzen von \(\lambda\) und \(r\),
\item für \(\log b\) nach den Potenzen von \(\beta,\lambda\) und \(r\),
\end{list}
